\textbf{COLUNI 2020 -- Questão 04}\par\medskip Leia o poema.

\par

[...] A cidade continuava lá, maciça, a de sempre, da sua infância, com suas ruas empoeiradas e seu casario humilde, embora sobre ela, membrana transparente, ao menos aos seus olhos, havia \underline{outra}, \underline{que} não mais correspondia àquela de antes. E se ele voltara ali em outras ocasiões, sem notar essa \underline{segunda}, dessa vez \underline{a} percebia não apenas num ou noutro detalhe, mas nela toda, cidade que se repovoava em sua memória.

\par

CARRASCOZA, J. A. \textit{Aos 7 e aos 40}. Rio de Janeiro: Alfaguara, 2016. p. 101-102.

\par

Os termos sublinhados no excerto acima se referem, respectivamente, a:\begin{enumerate}
\item outra; A cidade; outra; A cidade.
\item A cidade; outra; outra; outra.
\item outra; outra; A cidade; outra.
\item A cidade; outra; A cidade; outra.\end{enumerate}
